% Candidate article section. This file is a fragment, not a standalone document.
% Uses amsmath, amssymb, amsthm/the existing theorem and proof environments.
% All labels are prefixed ta: to avoid collisions with the submission.
% Read bounds.tex for the full derivations, qualifications, and literature audit.

\section{Finite-iteration bounds distinguishing Fanions}
\label{ta:section}

Identical worst-case exponents do not imply identical dependence on update
rank or learning rate. We first give a loss bound for the quadratic objective
used in our synthetic experiments. It is independent of the choice of
momentum and of the matrices supplied to the linear minimization oracle.

\begin{theorem}
\label{ta:obstruction}
Let $r=\min(m,n)$ and
\[
 f(X)=\frac12\|M^{1/2}(X-Z)N^{1/2}\|_{\mathrm F}^2,
 \qquad M,N\succeq0,
 \qquad L=\|M\|_{\mathrm{op}}\|N\|_{\mathrm{op}}.
\]
Put $E_0=M^{1/2}(X_0-Z)N^{1/2}$. Suppose
$X_{t+1}=X_t-\gamma_t A_t$, with $\gamma_t\geq0$,
$\operatorname{rank}(A_t)\leq k$, and $\|A_t\|_*\leq k$.
Then, for every $T$,
\begin{align}
 f(X_T)&\geq\frac12\sum_{j>kT}\sigma_j(E_0)^2,
 \label{ta:rank}\\
 f(X_T)&\geq\frac1{2r}
 \left[\|E_0\|_*-\sqrt L\,k\sum_{t=0}^{T-1}\gamma_t\right]_+^2.
 \label{ta:budget}
\end{align}
In particular, at constant learning rate $\gamma>0$, reaching
$f(X_T)\leq\epsilon$ requires
\[
 T\geq\frac{[\|E_0\|_*-\sqrt{2r\epsilon}]_+}
 {\sqrt L\,k\gamma}.
\]
\end{theorem}
\begin{proof}
Write $E_T=E_0-R_T$, where
$R_T=\sum_{t<T}\gamma_t M^{1/2}A_tN^{1/2}$ has rank at most $kT$.
The best rank-$kT$ approximation error of $E_0$ in Frobenius norm is
$\sum_{j>kT}\sigma_j(E_0)^2$, proving~\eqref{ta:rank}.
Also,
\[
 \|E_T\|_*\geq\|E_0\|_*-\sum_{t<T}\gamma_t
 \|M^{1/2}A_tN^{1/2}\|_*
 \geq\|E_0\|_*-\sqrt L\,k\sum_{t<T}\gamma_t.
\]
Combine this with $\|E_T\|_*\leq\sqrt r\|E_T\|_{\mathrm F}$
to obtain~\eqref{ta:budget}.
\end{proof}

An exact Fanion-$k$ update $A_t=\sum_{i=1}^k u_i v_i^\top$ satisfies both
constraints, regardless of the momentum matrix producing its singular
vectors. Consequently, at a shared learning rate, its accumulated nuclear
norm can grow at most $k\gamma$ per iteration. This provides a concrete
finite-horizon obstruction for low-rank Fanions. Increasing their learning
rates can reduce the second obstruction but cannot eliminate the first.
The additive update hypothesis is essential: normalization or decay of
$X_t$ must be accounted for separately.

The rank bound is sharp on $f(X)=\|X-I_r\|_{\mathrm F}^2/2$, $X_0=0$.
Every rank-$k$ method has relative loss at least
$[1-kT/r]_+$. Fanion-$k$ with exact line search and diagonal singular-vector
choices attains this bound, reaching zero after $\lceil r/k\rceil$
iterations. Muon solves this instance in one step.

F-Fanion updates have nuclear norm at most
$\alpha k+(1-\alpha)\sqrt r$; S-Fanion updates have nuclear norm at most
$\alpha k+(1-\alpha)\eta\sqrt{r mn}$. Thus~\eqref{ta:budget} extends with
these capacities in place of $k$. Their potentially full-rank updates
remove the original rank obstruction.

\subsection{Spectrum-dependent descent and its limitations}

For a true gradient $G$, write $S_k(G)=\sum_{i=1}^k\sigma_i(G)$ and
$D_k(G)=\sum_{i=1}^k u_i v_i^\top$. If $f$ is Frobenius $L$-smooth,
the step $\gamma=S_k(G)/(Lk)$ gives
\[
 f(X-\gamma D_k(G))\leq f(X)
 -\frac{S_k(G)^2}{2Lk}
 =f(X)-\frac{q_k(G)}{2L}\|G\|_{\mathrm F}^2,
 \qquad
 q_k(G)=\frac{S_k(G)^2}{k\|G\|_{\mathrm F}^2}.
\]
This follows directly from $\langle G,D_k\rangle=S_k(G)$ and
$\|D_k\|_{\mathrm F}^2\leq k$.
If $q_k(G_t)\geq\underline q_k>0$ along the trajectory, then
\[
 \min_{t<T}\|\nabla f(X_t)\|_{\mathrm F}^2
 \leq\frac{2L(f(X_0)-f_*)}{T\underline q_k}.
\]
Under a Frobenius PL inequality with constant $\mu$, the same calculation
gives loss contraction factor $1-\mu\underline q_k/L$.
For a flat singular spectrum, $q_k=k/r$, retaining an $r/k$ factor in
these bounds. At a common full-rank gradient,
\[
 \frac{q_r(G)}{q_1(G)}=
 \frac1r\left(\frac{\|G\|_*}{\|G\|_{\mathrm{op}}}\right)^2.
\]
Its value can be above or below one. These are comparisons of descent
certificates at common points, not ratios of observed training times.

For $a=\alpha$, $b=1-\alpha$, $x=S_k(G)/\|G\|_{\mathrm F}$, and
$\operatorname{rank}(G)\geq k$, the F-Fanion direction satisfies
\[
 q_{F,k}=\frac{(ax+b)^2}{a^2k+b^2+2abx}\geq q_k.
\]
Indeed, after multiplying by positive denominators, the difference has
numerator $(k-x^2)(2abx+b^2)\geq0$. For S-Fanions, Cauchy--Schwarz and
the triangle inequality give, when $a>0$,
\[
 q(G,D_{S,k})\geq
 \frac{q_k(G)}{\bigl(1+b\eta\sqrt{d_0}/(a\sqrt k)\bigr)^2},
 \qquad d_0=\#\{G_{ij}\ne0\}.
\]
These results quantify stability of full-rank hybrids. They also show why
Frobenius smoothness alone cannot establish a universal Muon advantage:
the normalized true gradient has cosine one.

To retain curvature, use a valid model
\[
 f(X-\gamma D)\leq f(X)-\gamma\langle G,D\rangle
 +\frac{\gamma^2}{2}C_X(D).
\]
The optimized decrease certificate is
$[\langle G,D\rangle]_+^2/(2C_X(D))$. For the quadratic,
$C_X(D)=\langle D,MDN\rangle$, and this certificate is exact.
Thus the appropriate structural comparison includes the singular spectrum,
curvature along the selected directions, and the actual oracle input.
No general norm convergence exponent by itself implies the experimentally
observed ordering across all tasks.
